\section{Casos y direcciones de investigación}

\subsection{Visual Agent Evidence}

La diapositiva clave muestra el acoplamiento entre la evidence matrix y el agent swarm:

\begin{figure}[H]
\centering
\includegraphics[width=0.8\textwidth]{slide-02.jpg}
\caption{La Slide 2 analiza la relación entre el agent swarm, la data quality y la governance, y adjunta una evidence matrix.}
\end{figure}

\subsection{Correspondencia entre las Slides y la práctica}

\begin{figure}[H]
\centering
\includegraphics[width=0.8\textwidth]{slide-03.jpg}
\caption{La Slide 3 muestra la integración del RLHF pipeline, el ops playbook y los checkpoint de gobernanza; es la evidencia visual de la importancia que la clase da a la actionability.}
\end{figure}

La Slide 3 enlaza la timeline, el monitoring y la governance en una matriz, y recuerda al equipo que en cada etapa debe haber un artifact entregable claramente definido.

\begin{knowledgebox}{Correspondencia entre las diapositivas y los entregables}
Antes del despliegue, convertir la timeline, los artifact y los checkpoints de la slide en un sprint backlog permite asegurar que cada deliverable tenga un responsable.
\end{knowledgebox}

\subsection{Research Directions}

La investigación futura puede centrarse en:

\begin{itemize}
    \item Conectar de forma automática (auto-connected) la reward model calibration con la policy evaluation
    \item La preference alignment con identidades Multi-agent
    \item Una Explainability layer que haga interpretables las restricciones KL
    \item Sim-to-real research para que el reward model se adapte a escenarios más amplios
\end{itemize}

Archit también insistió en la necesidad de cerrar el ciclo de retroalimentación entre la investigación teórica y la deployment real: las nuevas direcciones de investigación deben producir playbooks ejecutables, no solo paper-level insight.

\begin{knowledgebox}{Interpretabilidad de ingeniería de las direcciones de investigación}
Para que los resultados de investigación lleguen a la práctica, lo decisivo es la \textbf{interpretabilidad de ingeniería}: solo al explicar el KL margin, el reward model drift y los indicadores del agent dash board se consigue que los gerentes de producto y el equipo legal confíen en el alignment pipeline.
\end{knowledgebox}

\subsection{Resumen de la sección}

La parte de casos enlaza la evidencia visual de las slides, la estructura del agent swarm y las direcciones de investigación futura, y nos recuerda que el RLHF sigue evolucionando con rapidez.
